This chapter provides complementary information about the graphical user
interface, as seen by an \Eva{} user. Most of the documentation is available
in the graphical interface itself via the menu
\emph{Help $\rightarrow$ Documentation}.

\section{A general view of the GUI}

The \FramaC graphical user interface is currently best suited as a
{\em viewer} for analysis results. The typical use case is:

\begin{enumerate}
\item Run \FramaC in the command line and save the result
with option \texttt{-save proj.sav} to a session file;
\item Open the GUI with \texttt{frama-c-gui -load proj.sav} to load the saved session and
explore the analysis results interactively.
\end{enumerate}

The \FramaC GUI is not an IDE (e.g.~you cannot edit code with it). It is nevertheless very useful for academic examples and tutorials,
in which case it can obviate the need for the command line.

For large and industrial analyses, we suggest considering the GUI as a
complement, but not a full replacement, for the command-line interface.
Using a Makefile or script to parametrize and run the analysis ensures
reproducibility and helps iterative development.

\section{Eva-Related Views, Components, and Sidebar}

The \FramaC GUI is very rich and serves several other plug-ins and analyses.
To this end, it provides a wide range of components and a set of configurable
views; each {\em view} is a predefined arrangement of up to four components.

The top of the \emph{View and Components} sidebar lists available views;
clicking on a view opens it in a new tab.
In the active view, specific documentations can be accessed via the
\emph{?} button in the top right corner of each component.

\subsection{Views related to Eva}

The following predefined views are either related directly to \Eva{},
or often useful for navigation and code comprehension of \Eva{}-related properties:

\begin{description}
\item[Source Code]: this is the main view intended to navigate in the analyzed
source code, with various synthetic information on any variable, expression,
type or statement in the \emph{Inspector} panel.
\item[Properties]: not directly \Eva{}-related, but useful to browse the list
of alarms found by \Eva{}, and see which user assertions have been proved
(or not) by the analysis.
\item[Eva Summary]: useful first stop for getting an overview of the
analysis. It includes a textual summary, coverage information, profiling
information (Eva Flamegraph) and a list of all messages emitted during the analysis
(by default, only warnings and errors are shown).
\item[Eva Values]: the bread and butter of inspecting the set of possible values
inferred by \Eva{} for any variable or expression at any program point. It can
be very useful to understand the root causes of emitted alarms.
\item[Dive Dataflow]: the Dive plug-in presents graphs of dependencies
between memory locations, with various information from the \Eva{} analysis,
such as values inferred for these locations. It can be used to find the root
causes of an alarm by following the values of all locations involved from the
alarm point.
\item[Eva Mthread]: for inspection of concurrent properties after an analysis
with option \texttt{-mthread}.
\end{description}

\subsection{Sidebars}

\Eva{} has two dedicated sidebars in the \FramaC GUI.

The first one allows running, stopping and
restarting an analysis directly from the GUI. The most common analysis
parameters can be modified from this sidebar. As mentioned before, this is
especially useful for simple tests and quick comparisons, but for large,
longer analyses, we recommend using a Makefile or a script to ensure
reproducibility.

The second one allows selecting:
\begin{itemize}
\item a set of taints if the {\em taint} abstract domain (Section~\ref{sec:taint})
  was enabled. Only the selected taints (if any) are visible in the various
  components that show them (AST, Inspector, Properties…).
\item a set of callstacks, which restricts all Eva results displayed in the
  interface according to your selection — as if only these callstacks have been
  reached by the analysis. This allows focusing on some specific callstacks
  when inspecting the results of an analysis.
\end{itemize}
